\documentclass[a4paper,11pt, oneside,titlepage]{article}
\usepackage[pdftex]{graphicx}
\usepackage[latin2]{inputenc}
\usepackage{graphics}
\usepackage{amsmath,amssymb}
\pagestyle{plain}

\setlength{\topmargin}{-30pt}
\setlength{\headheight}{12pt}
\setlength{\headsep}{12pt}
\setlength{\oddsidemargin}{-17pt}
\setlength{\textheight}{57.0pc}
\setlength{\textwidth}{40.0pc}
 

\begin{document}
%%%%%%%%%%%%%%%%%%%%%%%%%%%%%%%%%%
%%%%%%%%%%%%%%%%%%%%%%%%%%%%%%%%%%

\pagestyle{empty} 


%\noindent Title: High-order fluctuations of temperature in hot QCD matter\\

%\noindent Authors: Jinhui Chen, Wei-jie Fu, Shi Yin, and Chunjian Zhang\\[0.3ex]


\noindent Dear editor,\\[-1ex]


\noindent We submit our manuscript titled ``High-order fluctuations of temperature in hot QCD matter" for consideration as a Letter in \textit{Physical Review Letters}. 
\\

\noindent Recent new measurements at the RHIC and LHC reveal that event-by-event mean transverse momentum fluctuations of charged particles, closely related to the temperature fluctuations, serve as an ideal probe of QCD thermodynamics. A new thermodynamic state function is introduced to describe the thermodynamics for the relevant heavy-ion collision experiments for the first time, allowing to compute the temperature fluctuations from the basic thermodynamic relations. The temperature fluctuations are suppressed remarkably as the system transitions from the hadron resonance gas (HRG) to the quark-gluon plasma (QGP) with increasing temperature or baryon chemical potential, alongside a negative skewness, in a general and model-independent way. These results provide a unique signature to discover
the thermodynamical temperature fluctuations in upcoming heavy-ion collision experiments and offer a novel approach to study the QCD phase diagram.\\

%\noindent Note that these findings are general and model-independent. They arises from the fact that the heat capacity of the QCD matter increases significantly from the HRG phase to the QGP phase.\\

\noindent We believe this research aligns with the high scientific standards of \textit{Physical Review Letters} and would greatly appreciate your consideration.\\




\noindent Sincerely yours, \\[0.3ex]

\noindent Jinhui Chen, Wei-jie Fu, Shi Yin, and Chunjian Zhang

\noindent 




%%%%%%%%%%%%%%%%%%%%%%%%%%%%%%%%%%
%%%%%%%%%%%%%%%%%%%%%%%%%%%%%%%%%%
\end{document}